\section{Perspectives}

Future work should first improve and evaluate the coarse diffusion model,
because it controls whether the speed advantage of WSGM can be obtained without
losing too much sample quality. The reference-coarse results suggest that the
conditional detail stages can work when the coarse state is correct; the next
step is therefore to make the generated coarse coefficient converge to the same
distribution as the reference coarse coefficients used during training.

It would then be useful to perform a systematic hyperparameter analysis. The
current results depend on several coupled choices: the wavelet family, the
decomposition level, the normalization of the coefficient channels, the number
of reverse diffusion steps, and the conditioning architecture. A controlled
study of these parameters would help determine whether the present quality loss
comes mainly from the representation, from the diffusion model, or from the
coarse-to-fine coupling.

Beyond sample generation, such models could also be used for downstream
physical tasks. One direction is rare-event prediction, where a fast generative
surrogate could be used to explore low-probability configurations without
running a large number of expensive DNS simulations. Another direction is
inverse-problem solving: given an observed density morphology, the model could
help search for initial conditions or coarse states that are likely to produce
that final shape. In this setting, the generative model would not only replace
forward simulations, but also provide a structured prior over physically
plausible RTI fields.
